\lesson{Topic SC-2 Use the Laws of Exponents}

\question{Question 1}

\blockhead{Question}
A person is designing a rectangular poster. The area of the poster is given as
the expression $3^{2} \times 3^{4}$ square inches. Use the laws of exponents to
simplify the expression and find the total area of the poster.

\checked{729 square inches}

\blockhead{What the question is asking}
This question asks you to find the number of square inches shown by
$3^{2} \times 3^{4}$.

\blockhead{Important words}
\begin{words}
  \item An \term{expression}is a math phrase made of numbers and operation signs.
  \item To \term{simplify}means rewrite an expression in a shorter form with the
    same value.
  \item A \term{factor}is a number being multiplied.
  \item The \term{base}is the number used as a repeated factor. In $3^{2}$, the
    base is 3.
  \item An \term{exponent}is the small raised number that tells how many times to
    use the base as a factor.
  \item A \term{power}is a base and its exponent together, such as $3^{2}$.
  \item A \term{product}is the answer to a multiplication problem.
  \item The \term{product-of-powers law}says that when powers have the same base,
    keep the base and add the exponents: $a^{m} \times a^{n} = a^{m+n}$.
  \item To \term{expand}a power means write all of its repeated factors.
  \item \term{Area}means the amount of flat space covered. It is measured here in
    square inches.
\end{words}

\blockhead{Work the problem one tiny step at a time}
\begin{steps}
  \item Look at the first power, $3^{2}$. Its base is 3.
  \item Look at the second power, $3^{4}$. Its base is also 3.
  \item The bases match, so the product-of-powers law applies.
  \item Keep the base 3.
  \item Add the exponents: $2 + 4 = 6$.
  \item Rewrite the expression: $3^{2} \times 3^{4} = 3^{6}$.
  \item Expand the power: $3^{6} = 3 \times 3 \times 3 \times 3 \times 3 \times 3$.
  \item Multiply the first two factors: $3 \times 3 = 9$.
  \item Multiply by the next 3: $9 \times 3 = 27$.
  \item Multiply by the next 3: $27 \times 3 = 81$.
  \item Multiply by the next 3: $81 \times 3 = 243$.
  \item Multiply by the last 3: $243 \times 3 = 729$.
\end{steps}

\finalans{729 square inches}

\blockhead{Why this makes sense}
The two powers contain two factors of 3 and four more factors of 3. That makes
six factors of 3 altogether. Expanding $3^{6}$ gives 729, so the poster area is
729 square inches.

\question{Question 2}

\blockhead{Question}
A person is baking cupcakes and makes a total of $2^{4}$ cupcakes. They want to
divide the batch equally among $2^{2}$ friends.
\par\noindent\runin{Part A:}What expression gives the number of cupcakes each
friend will get?
\par\noindent\runin{Part B:}How many cupcakes will each friend get?

\checked{Part A: $\dfrac{2^{4}}{2^{2}}$ (or $2^{4-2}$ or $2^{2}$); Part B: 4}

\blockhead{What the question is asking}
This question asks you to write the sharing problem and find how many cupcakes
one friend receives.

\blockhead{Important words}
\begin{words}
  \item An \term{operation}is a math action. Add joins amounts. Subtract takes
    away or finds a difference. Multiply makes equal groups. Divide splits into
    equal groups.
  \item An \term{expression}is a math phrase made of numbers and operation signs.
  \item \term{Equivalent}expressions look different but have the same value.
  \item To share \term{equally}means every friend gets the same amount.
  \item To \term{divide}means split a total into equal groups.
  \item The \term{dividend}is the total being divided. Here it is $2^{4}$
    cupcakes.
  \item The \term{divisor}is the number of equal groups. Here it is $2^{2}$
    friends.
  \item A \term{quotient}is the answer to a division problem.
  \item A \term{factor}is a number being multiplied.
  \item The \term{base}is the number used as a repeated factor. The base is 2 in
    both powers.
  \item An \term{exponent}is the small raised number that tells how many times to
    use the base as a factor.
  \item A \term{power}is a base and its exponent together, such as $2^{4}$.
  \item \term{Nonzero}means not equal to zero.
  \item The \term{quotient-of-powers law}says that when powers have the same
    nonzero base, keep the base and subtract the bottom exponent from the top
    exponent: $\dfrac{a^{m}}{a^{n}} = a^{m-n}$.
\end{words}

\blockhead{Work the problem one tiny step at a time}
\begin{steps}
  \item The total number of cupcakes is $2^{4}$.
  \item The number of friends is $2^{2}$.
  \item Sharing equally means divide the total by the number of friends.
  \item Write Part A: $\dfrac{2^{4}}{2^{2}}$.
  \item Both powers have base 2, so use the quotient-of-powers law.
  \item Subtract the exponents: $4 - 2 = 2$.
  \item Rewrite the quotient: $\dfrac{2^{4}}{2^{2}} = 2^{2}$.
  \item Expand $2^{2}$: $2^{2} = 2 \times 2$.
  \item Multiply: $2 \times 2 = 4$.
  \item Check by finding both powers: $2^{4} = 2 \times 2 \times 2 \times 2 = 16$
    and $2^{2} = 4$.
  \item Divide to check: $16 \div 4 = 4$.
\end{steps}

\finalans{Part A: $\dfrac{2^{4}}{2^{2}}$ (or an equivalent expression); Part B: 4
cupcakes per friend}

\blockhead{Why this makes sense}
Sixteen cupcakes split among four friends gives four cupcakes to each friend.
The exponent law and ordinary division agree.

\question{Question 3}

\blockhead{Question}
Evaluate using the power-of-a-power law: $(4^{3})^{2} = $ \blank[1.8cm].

\checked{4096}

\blockhead{What the question is asking}
This question asks you to find the number represented by $(4^{3})^{2}$.

\blockhead{Important words}
\begin{words}
  \item To \term{evaluate}means find the number value.
  \item \term{Parentheses}are curved marks that group math together.
  \item A \term{factor}is a number being multiplied.
  \item The \term{base}is the number used as a repeated factor. Here the base
    is 4.
  \item An \term{exponent}is the small raised number that tells how many times to
    use the base as a factor.
  \item A \term{power}is a base and its exponent together, such as $4^{3}$.
  \item A \term{power of a power}has one exponent on a power and another exponent
    outside the parentheses.
  \item The \term{power-of-a-power law}says to keep the base and multiply the
    exponents: $(a^{m})^{n} = a^{m \times n}$.
  \item To \term{expand}a power means write all of its repeated factors.
\end{words}

\blockhead{Work the problem one tiny step at a time}
\begin{steps}
  \item In $(4^{3})^{2}$, the base is 4.
  \item The exponent inside the parentheses is 3.
  \item The exponent outside the parentheses is 2.
  \item Use the power-of-a-power law.
  \item Multiply the exponents: $3 \times 2 = 6$.
  \item Keep the base 4: $(4^{3})^{2} = 4^{6}$.
  \item Expand: $4^{6} = 4 \times 4 \times 4 \times 4 \times 4 \times 4$.
  \item Multiply the first two factors: $4 \times 4 = 16$.
  \item Multiply by the next 4: $16 \times 4 = 64$.
  \item Multiply by the next 4: $64 \times 4 = 256$.
  \item Multiply by the next 4: $256 \times 4 = 1024$.
  \item Multiply by the last 4: $1024 \times 4 = 4096$.
  \item Check another way: $4^{3} = 64$, and $64^{2} = 64 \times 64 = 4096$.
\end{steps}

\finalans{4096}

\blockhead{Why this makes sense}
The outside exponent 2 means use the whole power $4^{3}$ twice. That makes six
factors of 4, so multiplying the exponents gives the same value as expanding.

\question{Question 4}

\blockhead{Question}
Evaluate: $7^{0} = $ \blank[1.8cm].

\checked{1}

\blockhead{What the question is asking}
This question asks for the number represented by 7 with a raised 0.

\blockhead{Important words}
\begin{words}
  \item To \term{evaluate}means find the number value.
  \item The \term{base}is the number used as a repeated factor. Here the base
    is 7.
  \item An \term{exponent}is the small raised number that tells how many times to
    use the base as a factor. Here the exponent is 0.
  \item A \term{power}is a base and its exponent together, such as $7^{0}$.
  \item \term{Nonzero}means not equal to zero.
  \item The \term{zero-exponent law}says any nonzero base raised to the power 0
    equals 1: $a^{0} = 1$ when $a \neq 0$.
\end{words}

\blockhead{Work the problem one tiny step at a time}
\begin{steps}
  \item Check the base: the base is 7.
  \item Seven is not zero, so the zero-exponent law applies.
  \item Follow the power pattern to see why.
  \item Start with $7^{2} = 7 \times 7 = 49$.
  \item Lower the exponent by 1 and divide by 7: $49 \div 7 = 7$, so $7^{1} = 7$.
  \item Lower the exponent by 1 again and divide by 7: $7 \div 7 = 1$, so
    $7^{0} = 1$.
  \item Apply the law directly: $7^{0} = 1$.
\end{steps}

\finalans{1}

\blockhead{Why this makes sense}
The power pattern divides by 7 each time the exponent drops by 1. The step from
$7^{1}$ to $7^{0}$ is $7 \div 7 = 1$. A zero exponent does not make the answer
zero.

\question{Question 5}

\blockhead{Question}
Simplify $\dfrac{4^{5}}{4^{2}} \cdot 4^{0} \cdot 4^{-1}$.

\checked{16}

\blockhead{What the question is asking}
This question asks you to find the number represented by the powers of 4 in the
expression.

\blockhead{Important words}
\begin{words}
  \item An \term{operation}is a math action. Add joins amounts. Subtract takes
    away or finds a difference. Multiply makes equal groups. Divide splits into
    equal groups.
  \item An \term{expression}is a math phrase made of numbers and operation signs.
  \item To \term{simplify}means rewrite an expression in a shorter form with the
    same value.
  \item To \term{evaluate}means find the number value of an expression.
  \item A \term{factor}is a number being multiplied.
  \item The \term{base}is the number used as a repeated factor. Every power here
    has base 4.
  \item An \term{exponent}is the small raised number that tells how many times to
    use the base as a factor.
  \item A \term{power}is a base and its exponent together, such as $4^{5}$.
  \item A \term{quotient}is the answer to a division problem.
  \item \term{Nonzero}means not equal to zero.
  \item The \term{quotient-of-powers law}says to keep a matching nonzero base and
    subtract exponents: $\dfrac{a^{m}}{a^{n}} = a^{m-n}$.
  \item The \term{zero-exponent law}says $a^{0} = 1$ for every nonzero base.
  \item A \term{fraction}$\dfrac{a}{b}$ means $a$ equal parts out of $b$ equal
    parts, where $b$ is not zero.
  \item A \term{reciprocal}is made by flipping a nonzero fraction. The reciprocal
    of $\dfrac{4}{1}$ is $\dfrac{1}{4}$.
  \item A \term{positive exponent}is an exponent greater than zero. A
    \term{positive power}uses a positive exponent, such as $4^{1}$.
  \item A \term{negative exponent}tells you to write the reciprocal of the
    matching positive power: $a^{-n} = \dfrac{1}{a^{n}}$.
  \item A \term{product}is the answer to multiplication.
  \item \term{PEMDAS}gives the work order: Parentheses, Exponents, Multiply and
    Divide from left to right, then Add and Subtract from left to right.
\end{words}

\blockhead{Work the problem one tiny step at a time}
\begin{steps}
  \item PEMDAS has no parentheses to work here.
  \item Work with the powers before completing the multiplication and division.
  \item Start with the quotient $\dfrac{4^{5}}{4^{2}}$.
  \item The bases match, so subtract the exponents: $5 - 2 = 3$.
  \item Rewrite the quotient: $\dfrac{4^{5}}{4^{2}} = 4^{3}$.
  \item Use the zero-exponent law: $4^{0} = 1$.
  \item Rewrite the negative exponent: $4^{-1} = \dfrac{1}{4^{1}}$.
  \item Evaluate $4^{1} = 4$, so $4^{-1} = \dfrac{1}{4}$.
  \item The expression is now $4^{3} \cdot 1 \cdot \dfrac{1}{4}$.
  \item Expand $4^{3}$: $4 \times 4 \times 4 = 64$.
  \item Multiply by 1: $64 \cdot 1 = 64$.
  \item Multiply by one fourth: $64 \cdot \dfrac{1}{4} = \dfrac{64}{4}$.
  \item Divide: $64 \div 4 = 16$.
\end{steps}

\finalans{16}

\blockhead{Why this makes sense}
The quotient leaves three factors of 4. The negative exponent contributes a
factor of one fourth, which removes one of those factors. Two factors of 4
remain, and $4^{2} = 16$.

\question{Question 6}

\blockhead{Question}
Simplify $(2^{3})^{2} \cdot (3 \cdot 2)^{2} \div \left(\dfrac{12}{2}\right)^{2}$.

\checked{64}

\blockhead{What the question is asking}
This question asks you to find the number represented by the whole expression.

\blockhead{Important words}
\begin{words}
  \item An \term{operation}is a math action. Add joins amounts. Subtract takes
    away or finds a difference. Multiply makes equal groups. Divide splits into
    equal groups.
  \item An \term{expression}is a math phrase made of numbers and operation signs.
  \item To \term{simplify}means rewrite an expression in a shorter form with the
    same value.
  \item To \term{evaluate}means find the number value of an expression.
  \item \term{Parentheses}are curved marks that group work together.
  \item A \term{factor}is a number being multiplied.
  \item The \term{base}is the number used as a repeated factor.
  \item An \term{exponent}is the small raised number that tells how many times to
    use the base as a factor.
  \item A \term{power}is a base and its exponent together, such as $2^{3}$.
  \item To \term{expand}a power means write all of its repeated factors.
  \item \term{Nonzero}means not equal to zero.
  \item To \term{multiply}means combine equal groups. A \term{product}is the
    answer to multiplication.
  \item The \term{power-of-a-power law}says $(a^{m})^{n} = a^{m \times n}$: keep
    the base and multiply the two exponents.
  \item The \term{power-of-a-product law}says $(ab)^{n} = a^{n}b^{n}$: raise each
    factor in the product to the exponent.
  \item To \term{divide}means split into equal groups. A \term{quotient}is the
    answer to division. In $a/b$, $a$ is divided by $b$.
  \item The \term{power-of-a-quotient law}says
    $\left(\dfrac{a}{b}\right)^{n} = \dfrac{a^{n}}{b^{n}}$ when $b \neq 0$.
  \item The condition $b \neq 0$ says $b$ must be nonzero.
  \item \term{Equal rank}means two operations share one place in the work order.
    Multiply and divide have equal rank, so do them from left to right.
  \item To \term{cancel}matching nonzero factors means divide them to make 1.
  \item \term{PEMDAS}gives the order: Parentheses, Exponents, Multiply and Divide
    from left to right, then Add and Subtract from left to right.
\end{words}

\blockhead{Work the problem one tiny step at a time}
\begin{steps}
  \item PEMDAS letter P says work inside parentheses first when possible.
  \item Inside the second group, multiply: $3 \cdot 2 = 6$.
  \item Inside the third group, divide: $12 \div 2 = 6$.
  \item Rewrite the expression: $(2^{3})^{2} \cdot 6^{2} \div 6^{2}$.
  \item PEMDAS letter E says work with exponents next.
  \item Use power of a power: $(2^{3})^{2} = 2^{3 \times 2}$.
  \item Multiply the exponents: $3 \times 2 = 6$.
  \item Rewrite the power: $(2^{3})^{2} = 2^{6}$.
  \item Expand: $2^{6} = 2 \times 2 \times 2 \times 2 \times 2 \times 2$.
  \item Multiply the first two factors: $2 \times 2 = 4$.
  \item Multiply by the next 2: $4 \times 2 = 8$.
  \item Multiply by the next 2: $8 \times 2 = 16$.
  \item Multiply by the next 2: $16 \times 2 = 32$.
  \item Multiply by the last 2: $32 \times 2 = 64$.
  \item Expand the other power: $6^{2} = 6 \times 6$.
  \item Multiply: $6 \times 6 = 36$.
  \item The expression is now $64 \cdot 36 \div 36$.
  \item PEMDAS says multiplication and division have equal rank, so work from
    left to right.
  \item To multiply $64 \cdot 36$, break 36 into $30 + 6$.
  \item Start $64 \cdot 30$ by finding $64 \cdot 3$.
  \item Multiply the tens: $60 \cdot 3 = 180$.
  \item Multiply the ones: $4 \cdot 3 = 12$.
  \item Add: $180 + 12 = 192$.
  \item Multiply by 10: $192 \cdot 10 = 1920$.
  \item So $64 \cdot 30 = 1920$.
  \item Now find $64 \cdot 6$.
  \item Multiply the tens: $60 \cdot 6 = 360$.
  \item Multiply the ones: $4 \cdot 6 = 24$.
  \item Add: $360 + 24 = 384$.
  \item Add the two partial products: $1920 + 384 = 2304$.
  \item Now divide $2304 \div 36$. Break 2304 into $2160 + 144$.
  \item Multiply: $36 \cdot 6 = 216$.
  \item Multiply by 10: $36 \cdot 60 = 2160$.
  \item So $2160 \div 36 = 60$.
  \item Multiply: $36 \cdot 4 = 144$.
  \item So $144 \div 36 = 4$.
  \item Add the two quotients: $60 + 4 = 64$.
  \item Check the matching factors: $36 \div 36 = 1$.
  \item Multiply: $64 \cdot 1 = 64$.
\end{steps}

\finalans{64}

\blockhead{Why this makes sense}
The second power and the third power are both $6^{2} = 36$. Multiplying by 36
and then dividing by 36 leaves the first value unchanged. The first value is
$(2^{3})^{2} = 64$.

\question{Question 7}

\blockhead{Question}
Simplify $5^{0} + 3^{-2} + \dfrac{2^{6} \cdot 2^{3}}{2^{7}}$.

\checked{$\mixed{5}{1}{9}$}

\blockhead{What the question is asking}
This question asks you to find the number made by adding the three parts of the
expression.

\blockhead{Important words}
\begin{words}
  \item An \term{operation}is a math action. Add joins amounts. Subtract takes
    away or finds a difference. Multiply makes equal groups. Divide splits into
    equal groups.
  \item An \term{expression}is a math phrase made of numbers and operation signs.
  \item To \term{simplify}means rewrite an expression in a shorter form with the
    same value.
  \item To \term{evaluate}means find the number value of an expression.
  \item A \term{factor}is a number being multiplied.
  \item \term{Nonzero}means not equal to zero.
  \item A \term{fraction}has a top number and a nonzero bottom number and names
    equal parts of a whole.
  \item A \term{whole number}has no fractional part, such as 5.
  \item The \term{numerator}is the top number of a fraction. The
    \term{denominator}is its bottom number.
  \item The \term{base}is the number used as a repeated factor.
  \item An \term{exponent}is the small raised number that tells how many times to
    use the base as a factor.
  \item A \term{power}is a base and its exponent together, such as $5^{0}$.
  \item The \term{zero-exponent law}says $a^{0} = 1$ when $a \neq 0$.
  \item To \term{flip}a fraction means swap its numerator (top number) and
    denominator (bottom number).
  \item A \term{reciprocal}is the result of flipping a nonzero fraction. The
    reciprocal multiplies with the original fraction to make 1. For example, flip
    $\dfrac{9}{1}$ to get its reciprocal $\dfrac{1}{9}$.
  \item A \term{negative exponent}tells you to use the reciprocal:
    $a^{-n} = \dfrac{1}{a^{n}}$.
  \item To \term{multiply}means combine equal groups. A \term{product}is the
    answer to multiplication.
  \item The \term{product-of-powers law}says $a^{m} \cdot a^{n} = a^{m+n}$ for
    matching bases.
  \item To \term{divide}means split into equal groups. A \term{quotient}is the
    answer to division.
  \item The \term{quotient-of-powers law}says $\dfrac{a^{m}}{a^{n}} = a^{m-n}$
    for a matching nonzero base.
  \item A \term{proper fraction}has a top number smaller than its bottom number,
    such as $\dfrac{1}{9}$.
  \item A \term{mixed number}has a whole number and a proper fraction, such as
    $\mixed{5}{1}{9}$.
  \item \term{PEMDAS}means Parentheses, Exponents, Multiply and Divide from left
    to right, then Add and Subtract from left to right.
\end{words}

\blockhead{Work the problem one tiny step at a time}
\begin{steps}
  \item PEMDAS says to work with powers before the final addition.
  \item Use the zero-exponent law: $5^{0} = 1$.
  \item Rewrite the negative exponent: $3^{-2} = \dfrac{1}{3^{2}}$.
  \item Evaluate the positive power: $3^{2} = 3 \times 3 = 9$.
  \item Therefore $3^{-2} = \dfrac{1}{9}$.
  \item In the numerator of the fraction, multiply powers with base 2.
  \item Add the numerator exponents: $6 + 3 = 9$.
  \item Rewrite the numerator: $2^{6} \cdot 2^{3} = 2^{9}$.
  \item Divide powers with base 2: $\dfrac{2^{9}}{2^{7}} = 2^{9-7}$.
  \item Subtract the exponents: $9 - 7 = 2$.
  \item Evaluate: $2^{2} = 2 \times 2 = 4$.
  \item The full expression is now $1 + \dfrac{1}{9} + 4$.
  \item Add the whole numbers: $1 + 4 = 5$.
  \item Keep the fractional part: $5 + \dfrac{1}{9} = \mixed{5}{1}{9}$.
\end{steps}

\finalans{$\mixed{5}{1}{9}$}

\blockhead{Why this makes sense}
The three parts are 1, one ninth, and 4. The whole-number parts total 5, and the
small positive fraction adds one ninth more. That gives $\mixed{5}{1}{9}$.

\question{Question 8}

\blockhead{Question}
Simplify $\dfrac{7^{4}}{7^{2}} \cdot 7^{0}$.

\checked{49}

\blockhead{What the question is asking}
This question asks you to find the number represented by the powers of 7.

\blockhead{Important words}
\begin{words}
  \item \term{PEMDAS}gives the work order: \textbf{P} means Parentheses;
    \textbf{E} means Exponents; \textbf{M} and \textbf{D} mean Multiply and
    Divide from left to right; \textbf{A} and \textbf{S} mean Add and Subtract
    from left to right.
  \item An \term{operation}is a math action. Add joins amounts. Subtract takes
    away or finds a difference. Multiply makes equal groups. Divide splits into
    equal groups.
  \item An \term{expression}is a math phrase made of numbers and operation signs.
  \item To \term{simplify}means rewrite an expression in a shorter form with the
    same value.
  \item A \term{factor}is a number being multiplied.
  \item The \term{base}is the number used as a repeated factor. Each power here
    has base 7.
  \item An \term{exponent}is the small raised number that tells how many times to
    use the base as a factor.
  \item A \term{power}is a base and its exponent together, such as $7^{4}$.
  \item \term{Nonzero}means not equal to zero.
  \item The \term{quotient-of-powers law}says to keep a matching nonzero base and
    subtract exponents: $\dfrac{a^{m}}{a^{n}} = a^{m-n}$.
  \item The \term{zero-exponent law}says every nonzero base to the power 0
    equals 1.
  \item A \term{product}is the answer to multiplication.
  \item A \term{quotient}is the answer to division.
  \item To \term{expand}a power means write all of its repeated factors.
  \item To \term{cancel matching factors}means divide the same nonzero factor
    from the top and bottom, which leaves a factor of 1.
\end{words}

\blockhead{Work the problem one tiny step at a time}
\begin{steps}
  \item PEMDAS says E (Exponents) comes before multiplication or division.
    Evaluate every power first.
  \item Expand $7^{4}$: $7 \times 7 \times 7 \times 7$.
  \item Multiply the first pair: $7 \times 7 = 49$.
  \item Multiply the second pair: $7 \times 7 = 49$.
  \item Multiply the two results: $49 \times 49 = 2401$.
  \item Expand $7^{2}$: $7 \times 7$.
  \item Multiply: $7 \times 7 = 49$.
  \item Use the zero-exponent law: $7^{0} = 1$.
  \item Rewrite the expression after all exponent work: $2401 \div 49 \cdot 1$.
  \item PEMDAS gives M and D equal rank, so work from left to right. Division
    appears first.
  \item Divide: $2401 \div 49 = 49$ because $49 \times 49 = 2401$.
  \item Multiply: $49 \cdot 1 = 49$.
  \item Check with the quotient law: $7^{4} \div 7^{2} = 7^{4-2} = 7^{2} = 49$.
\end{steps}

\finalans{49}

\blockhead{Why this makes sense}
Dividing four factors of 7 by two factors of 7 leaves two factors of 7. The zero
power contributes 1, so it does not change the product.

\question{Question 9}

\blockhead{Question}
Simplify $(3^{2})^{3}$.

\checked{729}

\blockhead{What the question is asking}
This question asks you to find the number represented by a power raised to
another power.

\blockhead{Important words}
\begin{words}
  \item An \term{operation}is a math action. Add joins amounts. Subtract takes
    away or finds a difference. Multiply makes equal groups. Divide splits into
    equal groups.
  \item An \term{expression}is a math phrase made of numbers and operation signs.
  \item To \term{simplify}means rewrite an expression in a shorter form with the
    same value.
  \item A \term{factor}is a number being multiplied.
  \item The \term{base}is the number used as a repeated factor. Here the base
    is 3.
  \item An \term{exponent}is the small raised number that tells how many times to
    use the base as a factor.
  \item A \term{power}is a base and its exponent together, such as $3^{2}$.
  \item A \term{power of a power}has an exponent outside parentheses that already
    contain a power.
  \item The \term{power-of-a-power law}says to keep the base and multiply the
    exponents: $(a^{m})^{n} = a^{m \times n}$.
  \item To \term{expand}means write all repeated factors.
\end{words}

\blockhead{Work the problem one tiny step at a time}
\begin{steps}
  \item Find the inner exponent: it is 2.
  \item Find the outer exponent: it is 3.
  \item Use the power-of-a-power law.
  \item Multiply the exponents: $2 \times 3 = 6$.
  \item Keep the base 3: $(3^{2})^{3} = 3^{6}$.
  \item Expand: $3^{6} = 3 \times 3 \times 3 \times 3 \times 3 \times 3$.
  \item Multiply: $3 \times 3 = 9$.
  \item Multiply: $9 \times 3 = 27$.
  \item Multiply: $27 \times 3 = 81$.
  \item Multiply: $81 \times 3 = 243$.
  \item Multiply: $243 \times 3 = 729$.
  \item Check: $3^{2} = 9$, so $(3^{2})^{3} = 9^{3} = 9 \times 9 \times 9
    = 81 \times 9 = 729$.
\end{steps}

\finalans{729}

\blockhead{Why this makes sense}
The outer exponent says to use $3^{2}$ three times. That makes $2 + 2 + 2 = 6$
factors of 3. Multiplying the exponents is the shortcut for counting those six
factors.

\question{Question 10}

\blockhead{Question}
Simplify $(2 \cdot 5)^{3}$.

\checked{1000}

\blockhead{What the question is asking}
This question asks you to find the number made by multiplying 2 and 5, then
using that result three times as a factor.

\blockhead{Important words}
\begin{words}
  \item An \term{operation}is a math action. Add joins amounts. Subtract takes
    away or finds a difference. Multiply makes equal groups. Divide splits into
    equal groups.
  \item An \term{expression}is a math phrase made of numbers and operation signs.
  \item To \term{simplify}means rewrite an expression in a shorter form with the
    same value.
  \item To \term{evaluate}means find the number value of an expression.
  \item A \term{product}is the answer to multiplication.
  \item A \term{factor}is a number being multiplied.
  \item \term{Parentheses}are curved marks that group work together.
  \item The \term{base}is the number used as a repeated factor. Here the grouped
    base is $(2 \cdot 5)$.
  \item An \term{exponent}is the small raised number that tells how many times to
    use the base as a factor.
  \item A \term{power}is a base and its exponent together, such as
    $(2 \cdot 5)^{3}$.
  \item To \term{expand}a power means write all of its repeated factors.
  \item The \term{power-of-a-product law}says $(ab)^{n} = a^{n}b^{n}$.
  \item \term{PEMDAS}gives the order: Parentheses, Exponents, Multiply and Divide
    from left to right, then Add and Subtract from left to right.
\end{words}

\blockhead{Work the problem one tiny step at a time}
\begin{steps}
  \item PEMDAS letter P says to work inside the parentheses first.
  \item Multiply inside the parentheses: $2 \cdot 5 = 10$.
  \item Rewrite the expression: $(2 \cdot 5)^{3} = 10^{3}$.
  \item The exponent 3 means use 10 as a factor three times.
  \item Expand: $10^{3} = 10 \times 10 \times 10$.
  \item Multiply the first two factors: $10 \times 10 = 100$.
  \item Multiply by the last factor: $100 \times 10 = 1000$.
  \item Check with the power-of-a-product law:
    $(2 \cdot 5)^{3} = 2^{3} \cdot 5^{3}$.
  \item Evaluate: $2^{3} = 2 \times 2 \times 2 = 8$.
  \item Evaluate: $5^{3} = 5 \times 5 \times 5 = 125$.
  \item Multiply to check: $8 \times 125 = 1000$.
\end{steps}

\finalans{1000}

\blockhead{Why this makes sense}
The grouped base is $2 \cdot 5 = 10$. Using 10 three times gives 1000.
Distributing the exponent to both factors also gives $8 \cdot 125 = 1000$.

\question{Question 11}

\blockhead{Question}
Simplify $6^{-2} + \left(\dfrac{8}{4}\right)^{3}$.

\checked{$\mixed{8}{1}{36}$}

\blockhead{What the question is asking}
This question asks you to find the number made by adding the two parts of the
expression.

\blockhead{Important words}
\begin{words}
  \item An \term{operation}is a math action. Add joins amounts. Subtract takes
    away or finds a difference. Multiply makes equal groups. Divide splits into
    equal groups.
  \item An \term{expression}is a math phrase made of numbers and operation signs.
  \item To \term{simplify}means rewrite an expression in a shorter form with the
    same value.
  \item To \term{evaluate}means find the number value of an expression.
  \item \term{Nonzero}means not equal to zero.
  \item A \term{fraction}has a top number and a nonzero bottom number and names
    equal parts of a whole.
  \item A \term{whole number}has no fractional part, such as 8.
  \item The \term{base}is the number used as a repeated factor.
  \item An \term{exponent}is the small raised number that tells how many times to
    use the base as a factor.
  \item A \term{power}is a base and its exponent together, such as $6^{2}$.
  \item A \term{reciprocal}is found by flipping a nonzero fraction. The
    reciprocal of $\dfrac{36}{1}$ is $\dfrac{1}{36}$.
  \item A \term{negative exponent}tells you to use the reciprocal of the matching
    positive power: $a^{-n} = \dfrac{1}{a^{n}}$.
  \item A \term{quotient}is the answer to division.
  \item A \term{proper fraction}has a top number smaller than its bottom number,
    such as $\dfrac{1}{36}$.
  \item An \term{improper fraction}has a top number at least as large as its
    bottom number, such as $\dfrac{289}{36}$.
  \item A \term{mixed number}has a whole number and a proper fraction.
  \item \term{PEMDAS}gives the order: Parentheses, Exponents, Multiply and Divide
    from left to right, then Add and Subtract from left to right.
\end{words}

\blockhead{Work the problem one tiny step at a time}
\begin{steps}
  \item PEMDAS letter P says to work inside the parentheses first.
  \item Divide inside the parentheses: $8 \div 4 = 2$.
  \item Rewrite that power: $\left(\dfrac{8}{4}\right)^{3} = 2^{3}$.
  \item PEMDAS letter E says to work with exponents next.
  \item Rewrite the negative exponent: $6^{-2} = \dfrac{1}{6^{2}}$.
  \item Evaluate $6^{2}$: $6 \times 6 = 36$.
  \item Therefore $6^{-2} = \dfrac{1}{36}$.
  \item Evaluate the other power: $2^{3} = 2 \times 2 \times 2$.
  \item Multiply: $2 \times 2 = 4$.
  \item Multiply: $4 \times 2 = 8$.
  \item The expression is now $\dfrac{1}{36} + 8$.
  \item Write the whole number first: $8 + \dfrac{1}{36}$.
  \item Combine the whole number and proper fraction:
    $8 + \dfrac{1}{36} = \mixed{8}{1}{36}$.
  \item Check as an improper fraction: $8 = \dfrac{288}{36}$, and
    $\dfrac{288}{36} + \dfrac{1}{36} = \dfrac{289}{36} = \mixed{8}{1}{36}$.
\end{steps}

\finalans{$\mixed{8}{1}{36}$}

\blockhead{Why this makes sense}
The negative power is a small positive amount, one thirty-sixth. The other power
equals 8. Their sum must be a little more than 8, and $\mixed{8}{1}{36}$ is a
little more than 8.

\question{Question 12}

\blockhead{Question}
Simplify $3^{4} \cdot 3^{2}$.

\checked{729}

\blockhead{What the question is asking}
This question asks you to find the number represented by multiplying two powers
of 3.

\blockhead{Important words}
\begin{words}
  \item An \term{operation}is a math action. Add joins amounts. Subtract takes
    away or finds a difference. Multiply makes equal groups. Divide splits into
    equal groups.
  \item An \term{expression}is a math phrase made of numbers and operation signs.
  \item To \term{simplify}means rewrite an expression in a shorter form with the
    same value.
  \item A \term{factor}is a number being multiplied.
  \item The \term{base}is the number used as a repeated factor. Both powers have
    base 3.
  \item An \term{exponent}is the small raised number that tells how many times to
    use the base as a factor.
  \item A \term{power}is a base and its exponent together, such as $3^{4}$.
  \item A \term{product}is the answer to multiplication.
  \item The \term{product-of-powers law}says to keep a matching base and add the
    exponents: $a^{m} \cdot a^{n} = a^{m+n}$.
  \item To \term{expand}a power means write all of its repeated factors.
\end{words}

\blockhead{Work the problem one tiny step at a time}
\begin{steps}
  \item Look at $3^{4}$. Its base is 3.
  \item Look at $3^{2}$. Its base is also 3.
  \item The bases match, so use the product-of-powers law.
  \item Keep the base 3.
  \item Add the exponents: $4 + 2 = 6$.
  \item Rewrite: $3^{4} \cdot 3^{2} = 3^{6}$.
  \item Expand: $3^{6} = 3 \times 3 \times 3 \times 3 \times 3 \times 3$.
  \item Multiply: $3 \times 3 = 9$.
  \item Multiply: $9 \times 3 = 27$.
  \item Multiply: $27 \times 3 = 81$.
  \item Multiply: $81 \times 3 = 243$.
  \item Multiply: $243 \times 3 = 729$.
\end{steps}

\finalans{729}

\blockhead{Why this makes sense}
Four factors of 3 and two more factors of 3 make six factors of 3. Expanding the
six factors confirms that the value is 729.
